% Generated draft: requires mathematical and source review.
\documentclass[11pt]{article}
\usepackage[margin=25mm]{geometry}
\usepackage{fontspec}
\setmainfont{Latin Modern Roman}
\usepackage{amsmath,amssymb,mathtools,unicode-math}
\setmathfont{Latin Modern Math}
\usepackage{enumitem,needspace,xurl,hyperref}
\setlength{\parindent}{0pt}
\setlength{\parskip}{5pt}
\setlength{\emergencystretch}{3em}
\newcommand{\N}{\mathbb{N}}
\newcommand{\Z}{\mathbb{Z}}
\newcommand{\Q}{\mathbb{Q}}
\newcommand{\R}{\mathbb{R}}
\newcommand{\C}{\mathbb{C}}
\newcommand{\F}{\mathbb{F}}
\newcommand{\A}{\mathbb{A}}
\newcommand{\PP}{\mathbb{P}}
\newcommand{\E}{\mathbb{E}}
\newcommand{\eps}{\varepsilon}
\newcommand{\dd}{\,\mathrm{d}}
\begin{document}
\section*{UnsolvedMath ID: 30003056}
\textbf{Source problem code:} catalogue entry
\subsection{Definitions and mathematical statement}
Isogeometric analysis (IGA) uses spline-type representations of geometry and approximation spaces. In a Galerkin discretization, a discrete unknown \(u_h\) in a trial space \(V_h\) is required to satisfy
\[
 a(u_h,v_h)=\ell(v_h)\qquad\text{for every }v_h\in W_h,
\]
where \(W_h\) is a test space and \(a\) and \(\ell\) denote the forms defining the variational problem. This notation describes the setting only; the source does not specify these spaces or forms, or restrict the problem to a particular differential equation.

For terminology, a quadrature rule approximates an integral by weighted evaluations:
\[
 I(f)=\int_{\Omega}f\,\mathrm d\mu,
 \qquad
 Q(f)=\sum_{j=1}^{N}w_j f(x_j).
\]
Here the integration domain \(\Omega\), measure \(\mu\), integrand class, nodes \(x_j\), weights \(w_j\), and number of evaluations \(N\) are unspecified inputs or design choices, not additional hypotheses supplied by the record. In Galerkin IGA, quadrature is used to approximate integrals occurring in the discrete variational forms.

\textbf{Source research question.} ``Can a general and effective quadrature solution be developed for Galerkin-based isogeometric analysis methods?''

\textbf{Unresolved formulation qualifications.} The source does not define the scope of ``general'' or the criterion for ``effective.'' In particular, it gives no class of admissible geometries, dimensions, spline spaces, coefficients, variational problems, or integrands, and no required accuracy, consistency, stability, convergence, or computational-cost bounds. It also does not specify whether effectiveness concerns individual integrals or the resulting Galerkin solution. Thus the supplied question is a research objective; it does not determine a unique quantified existence statement. No universal guarantee or missing hypotheses are imposed in this draft.
\subsection{Short English statement}
Develop a general and effective way to approximate the integrals arising in Galerkin isogeometric analysis. The intended range of methods and the required meaning of effectiveness remain unspecified, so a precise theorem-level target requires further formulation.
\subsection{Sources}
\begin{itemize}
\item Local UnsolvedMath registry: supplied complete record with permanent ID \(30003056\), title \emph{General Quadrature for Isogeometric Galerkin Methods}, and exact source problem code \texttt{catalogue entry}; collection: \emph{Oberwolfach Reports Open Problems}. The record gives proposed year 2016 and names no individual proposer.
\item Supplied report citation: \emph{Mini-Workshop: Mathematical Foundations of Isogeometric Analysis} (2016), \emph{Oberwolfach Reports}, pp.~341--385, DOI 10.4171/owr/2016/8. \url{https://doi.org/10.4171/owr/2016/8}. The supplied background describes Oberwolfach Reports as documenting mathematical workshops and their open-problem discussions; no problem-specific page is supplied.
\item Supplied literature reference, labeled ``Galerkin IGA quadrature paper'': \url{https://www.sciencedirect.com/science/article/abs/pii/S0045782514003235}.
\item Supplied literature reference, labeled ``2024 consistent-integration paper'': \url{https://doi.org/10.1016/j.cma.2024.117395}.
\item Literature assessment dated exactly \texttt{2026-08-21}: the supplied curation labels the problem open and reports that special and optimized quadrature rules, including consistency conditions, exist, while the requested general effective quadrature solution remains unresolved. This is supplied dated curation derived from the OWR record, not an independently verified literature assessment. The source explicitly requests independent verification of time-sensitive status and cited claims.
\item The registry labels the statement \texttt{exact} and describes it as clear, self-contained, and publication-ready. Those are supplied editorial qualifications; they do not resolve the unspecified mathematical scope and performance criteria identified above.
\end{itemize}
\end{document}
